\section{Multiple qubits and entanglement}

\subsection{Tensor products}

Two classical bits have four possibilities; two qubits have four basis vectors:
\[
\ket{00},\ \ket{01},\ \ket{10},\ \ket{11}.
\]
Independent systems combine with the tensor product $\otimes$. If
$\ket A=[a_0,a_1]^T$ and $\ket B=[b_0,b_1]^T$, then
\[
\ket A\otimes\ket B=
\begin{bmatrix}
a_0b_0\\a_0b_1\\a_1b_0\\a_1b_1
\end{bmatrix}.
\]
For $n$ qubits, the simulator stores $2^n$ amplitudes:
\[
\ket{\Psi}=\sum_{x=0}^{2^n-1}a_x\ket{x},\qquad
\sum_x|a_x|^2=1.
\]
This exponential growth is why state-vector simulation becomes expensive as
qubit count increases.

\subsection{Entanglement}

Some multi-qubit states cannot be factored into independent one-qubit states.
The Bell state
\[
\ket{\Phi^+}=\frac{\ket{00}+\ket{11}}{\sqrt2}
\]
is entangled. Each qubit measured alone is equally likely to be 0 or 1, but
their outcomes are perfectly correlated. It can be prepared from $\ket{00}$:
\begin{lstlisting}
H -I A -O A
CNOT -I A -O B
\end{lstlisting}
The first gate creates a superposition; the controlled gate correlates B with
A. Entanglement is not a classical alias. It is a property of the joint
amplitude vector.

\subsection{Controlled operations}

A controlled gate applies a target operation only on basis components where
its control condition is true. CNOT applies X when its control is one. If the
control is in superposition, linearity applies the operation to one branch and
not the other, which can create entanglement:
\[
\frac{\ket0+\ket1}{\sqrt2}\otimes\ket0
\xrightarrow{\mathrm{CNOT}}
\frac{\ket{00}+\ket{11}}{\sqrt2}.
\]
CCNOT requires both controls to be one. The system's OR and XOR derived gates
use ``any active'' and ``odd parity'' predicates respectively.

